\section{讨论}\label{sec:discussion}
\paragraph{机制而非格式本身。}AEIC的价值在于产生可缔约观测量，而非四个标签保证良好推理。Static AEIC在规则与观测量绑定时改变行为，Outcome Only说明一般奖励语言不足；但信号消融也不支持每个字段独立不可缺少的主张。

\paragraph{负迁移的重要性。}GLM结果表明，由Prompt介导的成本和信用通过后端特定策略解释，合同数值尚非通用货币。下一步应先在开发集估计各后端响应面，再在留出任务上测试固定合同；混合模型面板也应在解决校准后评估。

\paragraph{与拓扑和审议的关系。}生产激励决定是否产生信息性工作，通信拓扑决定谁能以何种成本观察这些工作。受控信息实验开始连接二者，却不把拓扑作为主因果变量。未来可在校准报告机制上优化信息流。

\paragraph{局限性。}主效应仅在一个主后端、三个二元答案QA域、一个Prompt行动菜单和多数聚合下建立。行动是LLM在指令下产生的行为，不代表类人偏好；信用是记账量而非货币；过程完整不等于事实验证。确认性实验检验固定静态规则，而非自适应Stackelberg最优性或一般均衡。更多调用和更长报告也带来真实计算成本。

\paragraph{更广泛含义。}该框架适用于中间工作可局部检查、执行器能够强制行动差异的场景，不应从结构完整但不可验证的报告制造虚假信心。因此记录合同、行动、执行路径和结算是贡献的一部分，而不只是工程日志。
